\title{DeepSeekMath: Advancing the Boundaries of Mathematical Reasoning in Open Language Models}

\begin{abstract}
Mathematical reasoning presents a considerable challenge for language models due to its intricate and structured characteristics. In this study, we present DeepSeekMath~7B, an extension of the pre-training process of DeepSeek-Coder-Base-v1.5 7B using 120B math-related tokens drawn from Common Crawl, along with natural language and code datasets. DeepSeekMath~7B attains an impressive 51.7\% accuracy on the competition-level MATH benchmark without utilizing external tools or ensemble methods, nearing the performance of Gemini-Ultra and GPT-4. Utilizing self-consistency with 64 samples from DeepSeekMath~7B results in a 60.9\% score on MATH. The enhanced mathematical reasoning capability of DeepSeekMath~is credited to two main aspects: first, leveraging the vast potential of publicly accessible web data through a carefully crafted data selection process; second, introducing Group Relative Policy Optimization (GRPO), a modified version of Proximal Policy Optimization (PPO) that improves mathematical reasoning performance while simultaneously optimizing PPO's memory efficiency.
\end{abstract}

\section{Introduction}

Large language models (LLMs) have transformed methods for mathematical reasoning within artificial intelligence, driving notable progress in both the quantitative reasoning benchmark \citep{MATH} and the geometry reasoning benchmark \citep{trinh2024solving}. Additionally, these models have been valuable in aiding humans to tackle challenging mathematical problems \citep{tao}. Nonetheless, state-of-the-art models like GPT-4 \citep{gpt4} and Gemini-Ultra \citep{gemini} remain inaccessible to the public, and currently available open-source models lag significantly behind in terms of performance.

In this research, we present DeepSeekMath, a specialized language model tailored for mathematics that substantially surpasses the mathematical performance of existing open-source models and nears the effectiveness of GPT-4 on academic evaluations. To accomplish this, we develop the DeepSeekMath Corpus, a large-scale, high-quality pre-training dataset containing 120B math tokens. This dataset is derived from the Common Crawl (CC) by utilizing a fastText-based classifier \citep{joulin2016fasttext}. Initially, the classifier is trained with positive samples from OpenWebMath \citep{openwebmath}, alongside a varied assortment of other web pages serving as negative samples. Next, the classifier is applied to extract additional positive examples from the CC, which are then further refined through manual annotation. The classifier is subsequently updated with this improved data to enhance its accuracy. Evaluation results demonstrate that the large-scale corpus is of excellent quality, as our foundational model DeepSeekMath-Base 7B attains scores of 64.2\% on GSM8K \citep{gsm8k} and 36.2\% on the competition-level MATH dataset \citep{MATH}, exceeding the performance of Minerva 540B \citep{minerva}. Furthermore, since the DeepSeekMath Corpus is multilingual, we observe gains on Chinese mathematical benchmarks \citep{wei2023cmath,agieval}. We view our methodology in mathematical data processing as an initial step for the research community, with ample opportunities for future enhancements.

DeepSeekMath-Base is initialized using DeepSeek-Coder-Base-v1.5 7B \citep{deepseek-coder}, as we find that beginning with a code-trained model is preferable to starting from a general LLM. Additionally, we observe that math-focused training enhances the model's performance on MMLU \citep{mmlu} and BBH benchmarks \citep{bbh}, suggesting that it not only improves the model's mathematical skills but also strengthens its overall reasoning abilities.

Following pre-training, we conduct mathematical instruction tuning on DeepSeekMath-Base using chain-of-thought \citep{cot}, program-of-thought \citep{pot,pal}, and tool-integrated reasoning \citep{tora} datasets. The resulting model, DeepSeekMath-Instruct 7B, surpasses all other 7B models and rivals 70B open-source instruction-tuned models.

Moreover, we propose Group Relative Policy Optimization (GRPO), a reinforcement learning (RL) variant of Proximal Policy Optimization (PPO) \citep{schulman2017proximal}. GRPO eliminates the need for a critic model by estimating the baseline from group scores, significantly reducing computational demands. Using only a subset of English instruction tuning data, GRPO achieves substantial gains over the strong DeepSeekMath-Instruct model, improving performance on both in-domain tasks (GSM8K: 82.9\% $\rightarrow$ 88.2\%, MATH: 46.8\% $\rightarrow$ 51.7\%) and out-of-domain mathematical challenges (e.g., CMATH: 84.6\% $\rightarrow$ 88.8\%) during RL training. We also present a unified framework to interpret various approaches such as Rejection Sampling Fine-Tuning (RFT) \citep{yuan2023scaling}, Direct Preference Optimization (DPO) \citep{dpo}, PPO, and GRPO. Within this framework, we recognize that all these methods can be viewed as either direct or simplified RL strategies. We conduct comprehensive experiments—including comparisons of online versus offline training, outcome-based versus process-based supervision, and single-turn versus iterative RL—to thoroughly explore the critical components of this paradigm. Finally, we clarify the reasons behind the performance improvements yielded by our RL approach on instruction-tuned models and propose future directions to develop more effective RL methods grounded in this unified framework.

\subsection{Contributions}
Our work presents scalable mathematical pre-training, along with an investigation and evaluation of reinforcement learning techniques.

\textbf{Math Pre-Training at Scale}

\begin{itemize}[topsep=0pt]
    \item We provide strong evidence that the publicly available Common Crawl dataset contains significant mathematical content. By employing a carefully crafted data selection pipeline, we create the DeepSeekMath Corpus, a high-quality dataset consisting of 120B tokens extracted from web pages filtered for mathematical relevance. This corpus is nearly 7 times larger than the math web pages dataset used by Minerva \citep{minerva} and 9 times larger than the recently introduced OpenWebMath \citep{openwebmath}.

    \item Our base model, DeepSeekMath-Base 7B, which is pre-trained on this dataset, achieves performance on par with Minerva 540B \citep{minerva}, demonstrating that the sheer number of parameters is not the sole determinant of mathematical reasoning capabilities. Smaller models pre-trained on high-quality data can also deliver strong results.

    \item We report insights from our experiments on math training. Performing code training before math training enhances the models’ proficiency in solving mathematical problems, both with and without the use of tools. This provides a partial resolution to the long-debated question: \textit{does code training enhance reasoning skills?} Our findings support that it does, at least in the context of mathematical reasoning.

    \item While training on arXiv papers is a common approach, particularly in many mathematics-focused studies, it does not yield significant improvements across the mathematical benchmarks evaluated in this paper.
\end{itemize}

\textbf{Exploration and Analysis of Reinforcement Learning}

\begin{itemize}[topsep=0pt]
    \item We propose Group Relative Policy Optimization (GRPO), a reinforcement learning algorithm that is both efficient and effective. GRPO eliminates the need for a critic model by estimating the baseline through group scores, which considerably lowers the training resource requirements compared to Proximal Policy Optimization (PPO).
    \item We show that GRPO notably improves the performance of our instruction-tuned model, DeepSeekMath-Instruct, relying solely on instruction-tuning data. Additionally, we observe improvements in out-of-domain generalization during the reinforcement learning phase.
    \item We offer a unified framework to interpret various methods including RFT, DPO, PPO, and GRPO. We also conduct extensive experiments—such as comparing online versus offline training, outcome versus process supervision, and single-turn versus iterative reinforcement learning—to thoroughly examine the core components of this framework.
    \item Building on our unified framework, we investigate the underlying factors contributing to the success of reinforcement learning and outline several promising directions for achieving more effective reinforcement learning for large language models (LLMs).
\end{itemize}

\subsection{Summary of Evaluations and Metrics}

\begin{itemize}[topsep=0pt]
    \item \textbf{English and Chinese Mathematical Reasoning}: We perform extensive evaluations of our models on English and Chinese benchmark datasets, spanning mathematical problems from elementary school to university level. The English benchmarks include GSM8K \citep{gsm8k}, MATH \citep{MATH}, SAT \citep{llemma}, OCW Courses \citep{minerva}, and MMLU-STEM \citep{mmlu}. The Chinese benchmarks comprise MGSM-zh \citep{mgsm}, CMATH \citep{wei2023cmath}, Gaokao-MathCloze \citep{agieval}, and Gaokao-MathQA \citep{agieval}. We assess the models’ capabilities to produce self-contained textual solutions without external tool use, as well as their proficiency in solving problems using Python.
    
    On English benchmarks, DeepSeekMath-Base is competitive with the proprietary Minerva 540B \citep{minerva} and outperforms all open-source base models (such as Mistral 7B \citep{mistral} and Llemma-34B \citep{llemma}), regardless of whether they have undergone math pre-training, often by a substantial margin. Importantly, DeepSeekMath-Base exhibits superior results on Chinese benchmarks, likely because we diverge from earlier approaches \citep{minerva,llemma} that collected exclusively English math pre-training data, incorporating high-quality non-English datasets instead. With mathematical instruction tuning and reinforcement learning, our resulting DeepSeekMath-Instruct and DeepSeekMath-RL models show robust performance, achieving an accuracy exceeding 50\% on the challenging MATH dataset for the first time within the open-source community.
\end{itemize}

\begin{itemize}
    \item \textbf{Formal Mathematics}: We assess DeepSeekMath-Base on the informal-to-formal theorem proving challenge introduced by \citep{dsp_proof}, using the miniF2F dataset \citep{minif2f} with Isabelle \citep{isabelle} as the selected proof assistant. DeepSeekMath-Base achieves notable few-shot autoformalization results.
    \item \textbf{Natural Language Understanding, Reasoning, and Code}: To establish a detailed evaluation of the models' broad comprehension, reasoning, and programming skills, we test DeepSeekMath-Base on the Massive Multitask Language Understanding (MMLU) benchmark \citep{mmlu}, which contains 57 multiple-choice tasks across a variety of domains; BIG-Bench Hard (BBH) \citep{bbh}, which includes 23 complex tasks primarily demanding multi-step reasoning; and also HumanEval \citep{codex} and MBPP \citep{mbpp}, both commonly used to assess code language models. Math pre-training enhances performance in both language understanding and reasoning.
\end{itemize}

\section{Math Pre-Training}

\subsection{Data Collection and Decontamination} This section describes the methodology used to assemble the DeepSeekMath Corpus from Common Crawl. As illustrated in Figure \ref{fig:data_collect}, we introduce an iterative pipeline outlining a systematic approach to gather a large-scale mathematical dataset from Common Crawl, beginning with a seed corpus (for example, a small yet high-quality math-related dataset). It should be emphasized that this method is extendable to other domains such as programming.

Initially, we select OpenWebMath \citep{openwebmath}, a high-quality compilation of mathematical web texts, as our seed corpus. From this corpus, we train a fastText model \citep{joulin2016fasttext} to identify additional mathematical web pages similar to those in OpenWebMath. Specifically, we randomly extract 500,000 samples from the seed corpus as positive examples, and 500,000 web pages from Common Crawl as negative examples. We utilize an open-source fastText implementation\footnote{\url{https://fasttext.cc}} with hyperparameters set to a 256-dimensional vector space, a learning rate of 0.1, word n-grams up to length 3, a minimum word count threshold of 3, and 3 training epochs. To reduce the scale of the raw Common Crawl data, we apply URL-based deduplication and near-deduplication, yielding 40 billion HTML pages. Using the trained fastText model, we retrieve mathematical web pages from this deduplicated set. To ensure quality, we score and rank the retrieved pages according to their fastText-predicted relevance and retain only the highest-ranked ones. The quantity of data retained is evaluated through pre-training trials on the top 40B, 80B, 120B, and 160B tokens. For the initial iteration, we choose to keep the top 40B tokens.

Following the initial round of data collection, a significant number of mathematical web pages remain uncollected, primarily because the fastText model is trained on a set of positive examples that lack sufficient diversity. To address this, we identify more mathematical web sources to expand the seed corpus, aiming to refine the fastText model. Specifically, we first partition the entire Common Crawl into separate domains; here, a domain is defined as web pages sharing the same base URL. For each domain, we compute the proportion of web pages collected during the first iteration. Domains where more than 10\% of web pages were gathered are classified as math-related (e.g., \url{mathoverflow.net}).

Next, we manually annotate URLs within these domains that correspond to mathematical content (e.g., \url{mathoverflow.net/questions}). Web pages linked to these URLs that were not previously collected are then added to the seed corpus. This strategy allows us to collect additional positive examples, enabling the training of an enhanced fastText model that can retrieve a greater amount of mathematical data in subsequent iterations. After four cycles of data collection, we accumulate 35.5M mathematical web pages, amounting to 120B tokens. In the fourth iteration, we observe that nearly 98\% of the data had already been captured by the third iteration, leading us to halt further data gathering.

To prevent contamination of benchmark datasets, we adopt the filtering approach from \cite{deepseek-coder}, removing web pages containing questions or answers from English mathematical benchmarks such as GSM8K \citep{gsm8k} and MATH \citep{MATH}, as well as Chinese benchmarks like CMATH \citep{wei2023cmath} and AGIEval \citep{agieval}. The filtering rule is as follows: any text segment that includes a 10-gram sequence exactly matching any substring from the evaluation benchmarks is excluded from our math training corpus. For benchmark texts shorter than 10 grams but at least 3 grams long, we apply exact matching to filter out contaminated pages.

\subsection{Evaluating the Quality of the DeepSeekMath~Corpus}

We conduct pre-training experiments to compare the DeepSeekMath~Corpus with recently released math-training corpora:

\begin{itemize}[topsep=0pt]
    \item \textbf{MathPile} \citep{mathpile}: a multi-source corpus (8.9B tokens) compiled from textbooks, Wikipedia, ProofWiki, CommonCrawl, StackExchange, and arXiv, with over 85\% of the data originating from arXiv;
    \item \textbf{OpenWebMath} \citep{openwebmath}: CommonCrawl data filtered for mathematical content, totaling 13.6B tokens;
    \item \textbf{Proof-Pile-2} \citep{llemma}: a mathematical corpus comprising OpenWebMath, AlgebraicStack (10.3B tokens of mathematical code), and arXiv papers (28.0B tokens). In experiments on Proof-Pile-2, we adhere to the arXiv:Web:Code data ratio of 2:4:1 as described in \cite{llemma}.
\end{itemize}

\subsubsection{Training Configuration}
\label{sec:quality-policy}
We perform math training on a general pre-trained language model comprising 1.3B parameters, utilizing the same architecture as the DeepSeek LLMs \citep{deepseek-llm}, referred to as DeepSeek-LLM 1.3B. Each mathematical dataset is used to train a separate model for a total of 150B tokens. All experiments employ the efficient and lightweight HAI-LLM \citep{haillm} training framework. Consistent with the training protocol of DeepSeek LLMs, we adopt the AdamW optimizer \citep{adamW} with parameters $\beta_1=0.9$, $\beta_2=0.95$, and $\mathrm{weight\_decay}=0.1$. The learning rate follows a multi-step schedule: it increases to its maximum over 2,000 warmup steps, then declines to 31.6\% of the peak at 80\% of training, and further drops to 10.0\% after 90\% of training completion. The peak learning rate is capped at 5.3e-4, with a batch size of 4 million tokens and a context window of 4K tokens.

\subsubsection{Results of Evaluation}

\textbf{The DeepSeekMath~Corpus is high-quality, multilingual, and the largest corpus available.}

\begin{itemize}[topsep=0pt]
    \item \textbf{High-quality}:
    We assess downstream task performance on eight mathematical benchmarks using few-shot chain-of-thought prompting \cite{cot}. As presented in Table \ref{tab:corpora_comparison}, the model trained on the DeepSeekMath~Corpus consistently outperforms others. Figure \ref{fig:corpora_comparisons} illustrates that the DeepSeekMath-trained model surpasses Proof-Pile-2 performance at 50B tokens (equivalent to one full epoch of Proof-Pile-2), suggesting that the average quality of the DeepSeekMath Corpus is superior.
    \item \textbf{Multilingual}:
    The DeepSeekMath~Corpus includes content in several languages, with English and Chinese being the most prominent. Table \ref{tab:corpora_comparison} shows that training on this corpus improves mathematical reasoning abilities in both English and Chinese. Conversely, existing corpora that focus mainly on English yield limited gains and may negatively impact performance in Chinese mathematical reasoning.
    \item \textbf{Large-scale}:
    The DeepSeekMath~Corpus is significantly larger than previously available mathematical corpora. As depicted in Figure \ref{fig:corpora_comparisons}, DeepSeek-LLM 1.3B trained on DeepSeekMath exhibits a steeper and more sustained learning progression. In comparison, baseline corpora are smaller, have been iterated over multiple times during training, and their corresponding model performances tend to plateau rapidly.
\end{itemize}

\subsection{Training and Evaluating DeepSeekMath-Base 7B}

In this section, we present DeepSeekMath-Base 7B, a foundational model with robust reasoning capabilities, particularly in mathematics. The model is initialized from DeepSeek-Coder-Base-v1.5 7B \citep{deepseek-coder} and trained on a dataset comprising 500 billion tokens. The data composition is as follows: 56\% originates from the DeepSeekMath Corpus, 4\% from AlgebraicStack, 10\% from arXiv, 20\% from Github code, and the remaining 10\% consists of natural language text from Common Crawl in both English and Chinese. Our training protocol largely follows the settings outlined in Section \ref{sec:quality-policy}, except that we employ a maximum learning rate of 4.2e-4 and utilize a batch size of 10 million tokens.

We perform an extensive evaluation of DeepSeekMath-Base 7B’s mathematical abilities, emphasizing its proficiency in generating self-contained mathematical solutions without external tool assistance, solving math problems with tools, and performing formal theorem proving. In addition to mathematical evaluation, we also offer a broader assessment of the base model’s capabilities, including natural language understanding, reasoning, and programming proficiency.

\paragraph{Mathematical Problem Solving with Step-by-Step Reasoning}

We assess DeepSeekMath-Base’s ability to solve mathematical problems using few-shot chain-of-thought prompting \citep{cot} across eight benchmarks in both English and Chinese. These benchmarks include quantitative reasoning datasets (such as GSM8K \citep{gsm8k}, MATH \citep{MATH}, and CMATH \citep{wei2023cmath}) as well as multiple-choice tasks (such as MMLU-STEM \citep{mmlu} and Gaokao-MathQA \citep{agieval}), covering a wide range of mathematical topics from basic to university-level difficulty.

As indicated in Table \ref{tab:base_math_cot}, DeepSeekMath-Base 7B achieves the highest performance across all eight benchmarks among open-source base models, which includes the popular general model Mistral 7B \citep{mistral} and the newly introduced Llemma 34B \citep{llemma}, the latter having undergone mathematical training on Proof-Pile-2 \citep{llemma}. Particularly, on the challenging competition-level MATH dataset, DeepSeekMath-Base outperforms all existing open-source base models by more than 10\% in absolute terms and even exceeds the performance of Minerva 540B \citep{minerva}, a proprietary base model 77 times larger that is based on PaLM \citep{palm} and further fine-tuned on mathematical literature.

\paragraph{Mathematical Problem Solving with Tool Use} We assess program-assisted mathematical reasoning on GSM8K and MATH datasets using few-shot program-of-thought prompting \citep{pot,pal}. The models are instructed to tackle each problem by crafting a Python program that can leverage libraries such as \textit{math} and \textit{sympy} for complex calculations. The output of the program's execution is then considered as the solution. As reported in Table \ref{tab:base_math_pot_proof}, DeepSeekMath-Base 7B surpasses the previous state-of-the-art Llemma 34B.

\paragraph{Formal Mathematics}

Formal proof automation plays a crucial role in guaranteeing the correctness and dependability of mathematical proofs while improving efficiency, an area attracting growing interest in recent years. We evaluate DeepSeekMath-Base 7B on the informal-to-formal proving task introduced in \citep{dsp_proof}, which involves generating a formal proof from an informal statement, its formal equivalent, and an informal proof. The evaluation is conducted on miniF2F \citep{minif2f}, a benchmark targeting formal Olympiad-level mathematics, where we produce formal proofs in Isabelle using few-shot prompting. Following the approach in \cite{dsp_proof}, models generate proof outlines, and the pre-built automated prover Sledgehammer \citep{sledgehammer} is used to complete missing steps. As demonstrated in Table \ref{tab:base_math_pot_proof}, DeepSeekMath-Base 7B shows robust capabilities in proof autoformalization.

\paragraph{Natural Language Understanding, Reasoning, and Code}

We measure the model’s natural language understanding on MMLU \citep{mmlu}, reasoning on BBH \citep{bbh}, and coding proficiency on HumanEval \citep{codex} and MBPP \citep{mbpp}. According to Table \ref{tab:understanding_reasoning_code}, DeepSeekMath-Base 7B exhibits marked improvements over its predecessor, DeepSeek-Coder-Base-v1.5 \citep{deepseek-coder}, on both MMLU and BBH, highlighting the beneficial effects of mathematical training on language comprehension and reasoning abilities. Moreover, by incorporating code tokens during continual training, DeepSeekMath-Base 7B successfully preserves the coding performance of DeepSeek-Coder-Base-v1.5 across both coding benchmarks. Overall, DeepSeekMath-Base 7B significantly outperforms the general-purpose model Mistral 7B \citep{mistral} on these three reasoning and coding evaluations.

\section{Supervised Fine-Tuning}

\subsection{SFT Data Preparation}

We develop a mathematical instruction-tuning dataset encompassing English and Chinese problems from various mathematical disciplines and varying degrees of difficulty. Each problem is accompanied by solutions presented in chain-of-thought (CoT) \citep{cot}, program-of-thought (PoT) \citep{pot,pal}, and tool-integrated reasoning formats \citep{tora}. The dataset contains a total of 776K training samples.

\begin{itemize}[topsep=0pt]
    \item \textbf{English mathematical datasets}: We annotate GSM8K and MATH problems with solutions that incorporate tool integration, and utilize a portion of MathInstruct \citep{MathInstruct} together with the training set of Lila-OOD \citep{lila}, where problems are addressed using CoT or PoT approaches. Our English dataset covers a broad range of mathematical areas including algebra, probability, number theory, calculus, and geometry.
    \item \textbf{Chinese mathematical datasets}: We gather Chinese K-12 math problems covering 76 sub-topics such as linear equations, with solutions provided both in CoT and tool-integrated reasoning formats.
\end{itemize}

\subsection{Training and Evaluation of DeepSeekMath-Instruct 7B}

Here, we present DeepSeekMath-Instruct 7B, which is fine-tuned on mathematical instructions starting from DeepSeekMath-Base. Training samples are randomly concatenated until reaching a maximum context length of 4K tokens. The model is trained for 500 steps using a batch size of 256 and a fixed learning rate of 5e-5.

We assess the models' mathematical capabilities both with and without tool usage across four quantitative reasoning benchmarks in English and Chinese. Our model is compared against leading contemporaneous models:

\begin{itemize}[topsep=0pt]
    \item \textbf{Closed-source models} include: (1) the GPT series, among which GPT-4 \citep{gpt4} and GPT-4 Code Interpreter~\footnote{\url{https://openai.com/blog/chatgpt-plugins\#\#code-interpreter}} are the most proficient, (2) Gemini Ultra and Pro \citep{gemini}, (3) Inflection-2 \citep{inflection-2}, (4) Grok-1~\footnote{\url{https://x.ai/model-card}}, as well as recently released models by Chinese companies, including (5) Baichuan-3~\footnote{\url{https://www.baichuan-ai.com}}, and (6) the newest GLM-4~\footnote{\url{https://open.bigmodel.cn/dev/api\#glm-4}} from the GLM family \citep{glm}.
\end{itemize}

These models serve general purposes, and the majority have been refined through a series of alignment processes.
\item \textbf{Open-source models} encompass: general models such as (1) DeepSeek-LLM-Chat 67B \citep{deepseek-llm}, (2) Qwen 72B \citep{qwen}, (3) SeaLLM-v2 7B \citep{seallm}, and (4) ChatGLM3 6B \citep{chatglm3}, along with models enhanced for mathematics including (5) InternLM2-Math 20B~\footnote{\url{https://github.com/InternLM/InternLM-Math}}, which is based on InternLM2 and underwent math-specific training followed by instruction tuning, (6) Math-Shepherd-Mistral 7B, which applies PPO training \citep{schulman2017proximal} to Mistral 7B \citep{mistral} using a process-supervised reward model, (7) the WizardMath series \citep{wizardmath}, which enhances mathematical reasoning capabilities in Mistral 7B and Llama-2 70B \citep{llama2} through evolve-instruct (an instruction tuning variant leveraging AI-evolved prompts) and PPO training with training tasks mainly drawn from GSM8K and MATH, (8) MetaMath 70B \citep{metamath}, a fine-tuned version of Llama-2 70B trained on an augmented GSM8K and MATH dataset, (9) ToRA 34B \cite{tora}, a CodeLlama 34B fine-tuned to incorporate tool-assisted mathematical reasoning, and (10) MAmmoTH 70B \citep{MathInstruct}, which is Llama-2 70B instruction-tuned on MathInstruct.
\end{itemize}

As indicated in Table \ref{tab:sft_rl_math}, when evaluated without the use of tools, DeepSeekMath-Instruct 7B exhibits strong step-by-step reasoning performance. Particularly, on the competition-level MATH dataset, our model outperforms all open-source models and most proprietary ones (e.g., Inflection-2 and Gemini Pro) by a margin of at least 9\% absolute. This holds true even compared to significantly larger models (e.g., Qwen 72B) or those specially enhanced with math-oriented reinforcement learning (e.g., WizardMath-v1.1 7B). While DeepSeekMath-Instruct is competitive with Chinese proprietary models such as GLM-4 and Baichuan-3 on MATH, it still trails behind GPT-4 and Gemini Ultra.

In the evaluation scenario where models can combine natural language reasoning with program-based tool usage for solving problems, DeepSeekMath-Instruct 7B achieves nearly 60\% accuracy on MATH, outperforming all current open-source models. On other benchmarks, it is on par with DeepSeek-LLM-Chat 67B, the previous state-of-the-art, despite being 10 times smaller.
\section{Reinforcement Learning}

\subsection{Group Relative Policy Optimization} Reinforcement learning (RL) has demonstrated effectiveness in further enhancing the mathematical reasoning abilities of large language models after the Supervised Fine-Tuning (SFT) phase \citep{wang2023math,wizardmath}. Here, we present our efficient and effective RL method, Group Relative Policy Optimization (GRPO).

\subsubsection{From PPO to GRPO} Proximal Policy Optimization (PPO) \citep{schulman2017proximal} is an actor-critic RL method frequently employed during RL fine-tuning of LLMs \citep{ouyang2022training}. Specifically, it improves LLMs by maximizing the following surrogate objective:

\begin{equation}
\footnotesize
    \mathcal{J}_{PPO}(\theta) = \mathbb{E}{[q \sim P(Q), o \sim \pi_{\theta_{old}}(O|q)]} \frac{1}{|o|} \sum_{t=1}^{|o|} \min \left[ \frac{\pi_\theta(o_{t} | q, o_{<t})}{\pi_{\theta_{old}}(o_{t} | q, o_{<t})} A_{t}, \text{clip} \left( \frac{\pi_\theta(o_{t} | q, o_{<t})}{\pi_{\theta_{old}}(o_{t} | q, o_{<t})}, 1 - \varepsilon, 1 + \varepsilon \right)  A_{t} \right]

where $\pi_{\theta}$ and $\pi_{\theta_{old}}$ denote the current and previous policy models, and $q, o$ represent questions and outputs sampled from the question dataset and the prior policy $\pi_{\theta_{old}}$, respectively. The parameter $\varepsilon$ is a clipping-related hyperparameter introduced in PPO to stabilize training. $A_t$ refers to the advantage, which is calculated using Generalized Advantage Estimation (GAE) \citep{gae}, based on the rewards $\{r_{\ge t}\}$ and a learned value function $V_{\psi}$. Therefore, in PPO, it is necessary to train a value function alongside the policy model. To prevent over-optimization of the reward model, the conventional method includes a per-token KL penalty from a reference model in the reward at each token \citep{ouyang2022training}, as expressed by

\begin{equation}
    r_{t} = r_\varphi(q, o_{\le t}) - \beta \log\frac{\pi_{\theta}(o_{t}|q, o_{<t})}{\pi_{ref}(o_{t}|q, o_{<t})},
\label{eq:PPO-reward}
\end{equation}

where $r_\varphi$ is the reward model, $\pi_{ref}$ is the reference model, typically the initial SFT model, and $\beta$ is the KL penalty coefficient.

Because the value function used in PPO is commonly another model of a similar size as the policy model, it imposes a significant computational and memory cost. Furthermore, during RL training, the value function acts as a baseline to reduce variance in advantage calculation. In the context of LLMs, since typically only the last token receives a reward from the reward model, training a value function that accurately estimates value at every token becomes challenging. To tackle this issue, as illustrated in Figure \ref{fig:grpo}, we introduce Group Relative Policy Optimization (GRPO), which eliminates the necessity for a separate value function approximation as in PPO. Instead, GRPO utilizes the average reward of multiple outputs sampled for the same question as the baseline. Concretely, for each question $q$, GRPO samples a group of outputs $\{o_1, o_2, \cdots, o_G\}$ from the old policy $\pi_{\theta_{old}}$ and optimizes the policy model by maximizing the following objective:

\begin{equation}
\footnotesize
\begin{split}
    \mathcal{J}_{GRPO}(\theta) &= \mathbb{E}{[q \sim P(Q), \{o_i\}_{i=1}^G \sim \pi_{\theta_{old}}(O|q)]}  \\
    & \frac{1}{G}\sum_{i=1}^G\frac{1}{|o_i|} \sum_{t=1}^{|o_i|} \left\{ \min \left[ \frac{\pi_\theta(o_{i,t} | q, o_{i,<t})}{\pi_{\theta_{old}}(o_{i,t} | q, o_{i,<t})} \hat{A}_{i,t}, \text{clip} \left( \frac{\pi_\theta(o_{i,t} | q, o_{i,<t})}{\pi_{\theta_{old}}(o_{i,t} | q, o_{i,<t})}, 1 - \varepsilon, 1 + \varepsilon \right)  \hat{A}_{i,t} \right] - \beta \mathbb{D}_{KL}\left[\pi_{\theta} || \pi_{ref}\right]\right\} ,
\end{split}
\label{eq:GRPO-obj}
\end{equation}

where $\varepsilon$ and $\beta$ are hyperparameters, and $\hat{A}_{i,t}$ denotes the advantage computed solely based on the relative rewards among outputs within each group, which will be elaborated on in the subsequent subsections. This group-relative method of calculating advantages in GRPO aligns naturally with the comparative framework of reward models, as these models are usually trained on datasets containing comparisons between outputs for the same question. Also, instead of incorporating the KL penalty within the reward, GRPO applies regularization by directly adding the KL divergence between the trained policy and the reference policy to the loss function, thus simplifying the calculation of $\hat{A}_{i,t}$. Unlike the KL penalty term in (\ref{eq:PPO-reward}), we approximate the KL divergence using the following unbiased estimator \citep{kl_approx}:

\begin{equation}
\small
    \mathbb{D}_{KL}\left

\begin{algorithm}[t]
  \small
  \caption{Iterative Group Relative Policy Optimization}
  \textbf{Input} initial policy model $\pi_{\theta_{\text{init}}}$; reward models $r_\varphi$; task prompts $\mathcal{D}$; 
  hyperparameters $\varepsilon$, $\beta$, $\mu$
  \begin{algorithmic}[1]
    \State Initialize policy model $\pi_\theta \leftarrow \pi_{\theta_{\text{init}}}$
    \For{iteration = 1, \dots, I}
       \State Set reference model $\pi_{ref} \leftarrow \pi_{\theta}$
      \For{step = 1, \dots, M}
      \State Draw a batch $\mathcal{D}_b$ from $\mathcal{D}$
      \State Save current policy $\pi_{\theta_{old}} \leftarrow \pi_{\theta}$ 
      \State For each question $q \in \mathcal{D}_b$, sample $G$ outputs $\{o_i\}_{i=1}^G \sim \pi_{\theta_{old}} (\cdot \mid q) $
      \State Evaluate rewards $\{r_i\}_{i=1}^{G}$ for each sampled output $o_i$ using $r_{\varphi}$ 
      \State Calculate $\hat{A}_{i,t}$ for the $t$-th token in $o_i$ by means of group relative advantage estimation.
      \For{GRPO iteration = 1, \dots, $\mu$}
        \State Update policy $\pi_{\theta}$ by maximizing the GRPO objective (Equation \ref{eq:GC-GRPO})
      \EndFor
    \EndFor \State Continuously train $r_\varphi$ using a replay mechanism. 
    \EndFor 
  \end{algorithmic}
  \textbf{Output} $\pi_\theta$
  \label{alg:iter-grpo}
\end{algorithm}

\subsubsection{Outcome Supervision RL with GRPO} Concretely, for each question $q$, a set of outputs $\{o_1, o_2, \cdots, o_G\}$ is generated from the previous policy $\pi_{\theta_{old}}$. These outputs are then scored by the reward model to obtain $G$ reward values $\mathbf{r}=\{r_1, r_2, \cdots, r_G\}$. These rewards are normalized by subtracting the group mean and dividing by the group standard deviation. Outcome supervision assigns the normalized reward as the advantage for all tokens in output $o_i$, i.e., $\hat{A}_{i, t} = \widetilde{r}_i = \frac{r_i- {\rm mean}(\mathbf{r})}{{\rm std}(\mathbf{r})}$, and the policy is optimized by maximizing the objective specified in equation (\ref{eq:GRPO-obj}).

\subsubsection{Process Supervision RL with GRPO} Outcome supervision delivers reward only after the entire output is generated, which may be inadequate or inefficient for guiding the policy in intricate mathematical problems. Following \cite{wang2023math}, we also investigate process supervision, where a reward is given at each reasoning step. Specifically, for question $q$ and sampled outputs $\{o_1, o_2, \cdots, o_G\}$, the process reward model assigns rewards to each step, resulting in rewards: $\mathbf{R} = \{ \{r_1^{index(1)},\cdots,r_1^{index(K_1)}\}, \cdots,  \{r_G^{index(1)},\cdots,r_G^{index(K_G)}\} \}$, where $index(j)$ is the ending token index of the $j$-th step, and $K_i$ denotes the number of steps in the $i$-th output. These rewards are normalized by the overall mean and standard deviation, i.e., $\widetilde{r}_i^{index(j)} = \frac{r_i^{index(j)} - {\rm mean}(\mathbf{R})}{{\rm std}(\mathbf{R})}$. Then, the advantage for each token is computed as the cumulative sum of normalized rewards from subsequent steps, i.e., $\hat{A}_{i, t} = \sum_{index(j) \ge t} \widetilde{r}_i^{index(j)}$, after which the policy is improved by maximizing the objective in equation (\ref{eq:GRPO-obj}).

\subsubsection{Iterative RL with GRPO} As the reinforcement learning training advances, the previous reward model may become inadequate to guide the current policy model effectively. Hence, we investigate an iterative RL approach using GRPO. As illustrated in Algorithm \ref{alg:iter-grpo}, iterative GRPO involves generating new training datasets for the reward model derived from samples produced by the policy model. We continuously update the existing reward model through a replay mechanism that includes 10\% of past data. Subsequently, the policy model is used as the reference model, and training continues using the updated reward model.

\subsection{Training and Evaluating DeepSeekMath-RL}

We perform reinforcement learning starting from DeepSeekMath-Instruct 7B. The RL training dataset consists of chain-of-thought style questions related to GSM8K and MATH taken from the SFT dataset, totaling approximately 144K questions. Other SFT questions are excluded to examine the effect of RL on benchmarks with limited data during the RL process. The reward model training set is constructed following the methodology in \citep{wang2023math}. Our initial reward model is trained based on DeepSeekMath-Base 7B with a learning rate of 2e-5. For GRPO, the policy model’s learning rate is set to 1e-6, and the KL coefficient is fixed at 0.04. For each question, 64 outputs are sampled. The maximum sequence length is 1024, and the batch size for training is 1024. The policy model undergoes only one update after each exploration phase. Evaluation of DeepSeekMath-RL 7B is performed on benchmarks in the same manner as DeepSeekMath-Instruct 7B. For DeepSeekMath-RL 7B, GSM8K and MATH with chain-of-thought reasoning are considered in-domain tasks, while the remaining benchmarks are treated as out-of-domain tasks.

Table \ref{tab:sft_rl_math} presents the results of both open- and closed-source models employing chain-of-thought and tool-augmented reasoning on English and Chinese benchmarks. Our observations include: 1) DeepSeekMath-RL 7B achieves accuracies of 88.2\% on GSM8K and 51.7\% on MATH when using chain-of-thought reasoning. This level of performance exceeds that of all open-source models within the 7B to 70B parameter range, as well as most closed-source models. 2) Importantly, DeepSeekMath-RL 7B is trained solely on chain-of-thought instruction tuning data from GSM8K and MATH, starting from the DeepSeekMath-Instruct 7B model. Despite the limited scope of its training data, it outperforms DeepSeekMath-Instruct 7B on every evaluation metric, highlighting the benefits of reinforcement learning.

\section{Discussion}
In this section, we present the insights gained from our pre-training and reinforcement learning experiments.

\subsection{Insights Gained from Pre-Training}

We begin by sharing our experiences related to pre-training. Unless otherwise stated, the training configurations follow those described in Section \ref{sec:quality-policy}. It should be noted that when we mention the DeepSeekMath Corpus in this section, we refer to the 89B-token dataset collected during the second iteration of data gathering.

\subsubsection{Code Training Enhances Mathematical Reasoning}

A widely held but unconfirmed assumption is that code training enhances reasoning capabilities. We seek to provide a partial answer to this claim, specifically within the mathematical domain: code training improves models’ capacity for mathematical reasoning both with and without the use of external tools.

To investigate the impact of code training on mathematical reasoning, we conducted experiments using two different approaches: two-stage training and one-stage training:

\textbf{Two-Stage Training}

\begin{itemize}[topsep=0pt]
    \item \textbf{Code Training for 400B Tokens $\rightarrow$ Math Training for 150B Tokens}: We first train DeepSeek-LLM 1.3B on 400 billion code tokens, followed by 150 billion math tokens;
    \item \textbf{General Training for 400B Tokens $\rightarrow$ Math Training for 150B Tokens}: As a baseline comparison, we also conduct training with general tokens (drawn from a large-scale general corpus developed by DeepSeek-AI) instead of code tokens in the initial training phase, to explore whether code tokens offer advantages over general tokens in enhancing mathematical reasoning.
\end{itemize}

\textbf{One-Stage Training}

\begin{itemize}[topsep=0pt]
    \item \textbf{Math Training for 150B Tokens}: We train DeepSeek-LLM 1.3B solely on 150 billion math tokens;
    \item \textbf{Training on a mixture of 400B Code Tokens and 150B Math Tokens}: Since math training following code training leads to a decline in coding performance, we examine whether blending code tokens with math tokens in a single-stage training can still improve mathematical reasoning while also reducing catastrophic forgetting.
\end{itemize}

\paragraph{Results}
Tables \ref{tab:code-math} and \ref{tab:code-math-general-eval} present the downstream task performances under various training configurations.

Training on code data enhances program-assisted mathematical reasoning, both in two-stage and one-stage training frameworks. As evidenced in Table \ref{tab:code-math}, within the two-stage setup, code training alone considerably improves performance on GSM8K and MATH tasks solved via Python. Subsequent math training further boosts these capabilities. Notably, in the one-stage scenario, mixing code and math tokens effectively alleviates the catastrophic forgetting observed in two-stage training, while simultaneously supporting both coding (Table \ref{tab:code-math-general-eval}) and program-assisted mathematical reasoning (Table \ref{tab:code-math}).

Moreover, code training enhances mathematical reasoning without reliance on tools. In the two-stage approach, the initial code training phase alone yields moderate improvements, which then accelerate the effectiveness of later math training, culminating in the highest performance levels. Conversely, combining code and math tokens for one-stage training diminishes performance on mathematical reasoning without tool use. This may be because DeepSeek-LLM 1.3B, given its relatively small scale, lacks sufficient capacity to effectively integrate both code and mathematical data simultaneously.

\subsubsection{ArXiv Papers Appear Ineffective for Enhancing Mathematical Reasoning}

ArXiv papers are frequently incorporated as part of math pre-training datasets \citep{minerva,gpt-f,llemma,mathpile}. Nonetheless, thorough investigations into their influence on mathematical reasoning remain limited. Interestingly, our experiments suggest that arXiv papers do not significantly aid in improving mathematical reasoning. We conduct experiments on models of varying scales, including DeepSeek-LLM 1.3B and DeepSeek-Coder-Base-v1.5 7B \citep{deepseek-coder}, utilizing arXiv corpora processed through different pipelines:

\begin{itemize}[topsep=0pt]
    \item \textbf{MathPile} \citep{mathpile}: an 8.9B-token dataset constructed with cleaning and filtering heuristics, containing over 85\% scientific arXiv papers;
    \item \textbf{ArXiv-RedPajama} \citep{redpajama}: the full collection of arXiv LaTeX files with preambles, comments, macros, and bibliographies stripped out, amounting to 28.0B tokens.
\end{itemize}

In our trials, DeepSeek-LLM 1.3B is trained for 150B tokens, and DeepSeek-Coder-Base-v1.5 7B for 40B tokens on each arXiv dataset. The results indicate that arXiv papers do not effectively enhance mathematical reasoning. When models are trained exclusively on arXiv data, they show no significant gains and sometimes even regress across various mathematical benchmarks of differing difficulty levels used in this work. These benchmarks encompass quantitative reasoning tasks like GSM8K and MATH (see Table \ref{tab:arxiv-cot}), multiple-choice tests such as MMLU-STEM (Table \ref{tab:arxiv-cot}), as well as formal mathematics evaluations like miniF2F (Table \ref{tab:arxiv_atp}).

Nevertheless, this finding has limitations and should be interpreted cautiously. We have not yet explored:

\begin{itemize}[topsep=0pt]
    \item The influence of arXiv tokens on particular math-related tasks outside the scope of this study, for example, theorem informalization, which involves converting formal statements or proofs into informal versions;
    \item The effect of arXiv tokens when combined with other data sources;
    \item Whether benefits from arXiv papers might emerge at larger model scales.
\end{itemize}

Hence, further investigation is needed, which we defer to future work.

\subsection{Insights of Reinforcement Learning}
\subsubsection{Towards a Unified Paradigm} In this section, we present a unified framework to analyze various training approaches such as SFT, RFT, DPO, PPO, and GRPO, and perform experiments to investigate elements of this unified framework. Generally, the gradient with respect to the parameter $\theta$ for a training method can be formulated as:

\begin{equation}
    \nabla_{\theta}\mathcal{J}_{\textcolor{red}{\mathcal{A}}}(\theta) = \mathbb{E}[\underbrace{(q,o) \sim \textcolor{red}{\mathcal{D}}}_{Data \ Source}]\left( \frac{1}{|o|} \sum_{t=1}^{|o|}  \underbrace{GC_{{\mathcal{A}}}(q, o, t, \textcolor{red}{\pi_{{rf}}})}_{Gradient \ Coefficient}  \nabla_{\theta}\log \pi_{\theta}(o_t | q, o_{<t})\right).
\label{eq:objective}
\end{equation}

This framework consists of three primary elements: 1) \textit{Data Source $\mathcal{D}$}, which specifies the dataset used for training; 2) \textit{Reward Function $\pi_{{rf}}$}, providing the reward signal during training; 3) \textit{Algorithm $\mathcal{A}$}, which transforms the training data and reward signal into the gradient coefficient $GC$ that controls the scale of the penalty or reward applied to the data. We examine several notable methods within this unified framework:

\begin{itemize}
    \item  \textbf{Supervised Fine-tuning (SFT)}: SFT adapts a pretrained model by fine-tuning it on a curated human-labeled dataset.
    \item \textbf{Rejection Sampling Fine-tuning (RFT)}: RFT continues fine-tuning the SFT model using filtered outputs generated from SFT model responses to SFT questions, where filtering is based on answer correctness.
    \item \textbf{Direct Preference Optimization (DPO)}: DPO further adjusts the SFT model by fine-tuning it on augmented outputs sampled from the SFT model, employing a pairwise DPO loss.
    \item \textbf{Online Rejection Sampling Fine-tuning (Online RFT)}: Unlike RFT, Online RFT starts with the SFT model as the policy model and refines it by fine-tuning on augmented outputs drawn from the current policy model in real time.
    \item \textbf{PPO/GRPO}: PPO/GRPO initialize the policy model using the SFT model and enhance it through reinforcement with outputs sampled from the ongoing policy model.
\end{itemize}

We provide a summary of the components of these methods in Table \ref{tab:data_policy}. For a more comprehensive derivation process, please consult Appendix \ref{app:analysis-rl}.

\paragraph{Observation about Data Source} We categorize the data source into two types: online sampling and offline sampling. Online sampling indicates that the training data comes from the exploration outcomes of the current training policy model, whereas offline sampling means the training data is derived from the sampling results of the initial SFT model. RFT and DPO adopt the offline approach, while Online RFT and GRPO utilize the online approach.

As illustrated in Figure \ref{fig:rl-analysis}, we observe that Online RFT markedly outperforms RFT across two benchmarks. More precisely, Online RFT performs similarly to RFT during the early phases of training but gains a clear advantage in the later stages, highlighting the benefits of online training. This makes intuitive sense since, at the beginning, the actor and the SFT model are quite similar, causing the sampled data to show only minor differences. However, as training progresses, data sampled from the actor reflects more substantial differences, making real-time data sampling more advantageous.

\paragraph{Observation about Gradient Coefficient} The algorithm modifies the gradient coefficient using the reward signal to update model parameters. In our experiments, we classify the reward function into two categories: `Rule' and `Model'. The `Rule' category assesses the response quality based on the correctness of the answer, while the `Model' category involves training a reward model to assign scores to each response. The reward model’s training data is grounded in the rule-based judgment. Equations \ref{eq:GC-RFT} and \ref{eq:GC-GRPO} underscore a fundamental distinction between GRPO and Online RFT: GRPO uniquely adjusts its gradient coefficient according to the reward values generated by the reward model, enabling differential reinforcement and punishment based on varying response quality. In contrast, Online RFT does not include this mechanism; it does not penalize incorrect responses and uniformly reinforces all correct answers with the same level of intensity.

As illustrated in Figure \ref{fig:rl-analysis}, GRPO outperforms online RFT, underscoring the effectiveness of modifying positive and negative gradient coefficients. Moreover, GRPO+PS exhibits better results than GRPO+OS, demonstrating the advantages of employing fine-grained, step-aware gradient coefficients. Additionally, we investigate iterative RL by conducting two rounds of iteration in our experiments. As depicted in Figure \ref{fig:iter-rl}, iterative RL markedly enhances performance, particularly after the first iteration.

\subsubsection{Why Does RL Work?} In this study, we apply reinforcement learning on a subset of instruction tuning data, achieving notable performance gains over the instruction tuning model. To further elucidate why reinforcement learning is effective, we assess the Pass@K and Maj@K accuracy of the Instruct and RL models on two benchmark datasets. As shown in Figure \ref{fig:maj-pass}, RL improves Maj@K performance but does not enhance Pass@K. These results suggest that RL improves the model’s overall robustness in output distribution; in other words, \textbf{the improvement appears to stem from increasing the likelihood of selecting the correct response within the TopK rather than from enhancing the model’s fundamental capabilities.} Similarly, \citep{wang2023making} uncovered a \textbf{misalignment problem} in reasoning tasks within the SFT model, demonstrating that reasoning performance can be elevated through preference alignment strategies \citep{yuan2023rrhf,song2023preference,wang2023making}.

\subsubsection{How Can We Achieve More Effective RL?}
\label{sec:effec_rl}
We show that RL performs exceptionally well in mathematical reasoning tasks and offer a unified framework to understand various representative training methods. Within this framework, all methods are framed as either direct or simplified RL approaches. As summarized in Equation \ref{eq:objective}, three crucial components exist: Data Source, Algorithm, and Reward Function. We outline some prospective directions regarding these components.

\paragraph{Data Source} The data source constitutes the fundamental input for all training methods. Specifically for RL, the data source refers to unlabeled questions with outputs sampled from the policy model. In this paper, we use only the questions from the instruction tuning phase and apply a basic nucleus sampling method to generate outputs. We believe this limitation partly explains why our RL pipeline solely enhances Maj@K performance. Future work will investigate the RL pipeline using out-of-distribution question prompts combined with \textbf{advanced sampling (decoding) strategies}, such as those based on tree-search algorithms \citep{yao2023tree}. Additionally, \textbf{efficient inference techniques} \citep{xia-etal-2023-speculative,leviathan2023fast,kwon2023efficient,xia2024unlocking}, which impact the exploration efficiency of policy models, will also be critically important.

\paragraph{Algorithms} Algorithms utilize the data and reward signals to adjust the gradient coefficient, thereby updating the model parameters. According to Equation \ref{eq:objective}, current methods to a large extent fully \textbf{TRUST} the reward function's signal to either increase or decrease the conditional probability of specific tokens. Nevertheless, guaranteeing the reward signal's reliability is not feasible, especially for highly complex tasks. For instance, the PRM800K datasets \citep{lightman2023let}, despite being carefully annotated by skilled annotators, still contain about 20\% incorrect annotations\footnote{\url{https://github.com/openai/prm800k/issues/12\#issuecomment-1728491852}}. In response, we aim to investigate reinforcement learning algorithms that maintain robustness in the presence of noisy reward signals. We anticipate that such \textbf{WEAK-TO-STRONG} \citep{burns2023weak} alignment methods will fundamentally transform learning algorithms.

\paragraph{Reward Function} The reward function provides the training signal in RL, typically implemented as a neural reward model. We identify three crucial directions for reward models: 1) \textbf{Improving the generalization capability of the reward model.} The reward model must generalize effectively to manage out-of-distribution queries and complex decoding outputs; otherwise, reinforcement learning may only serve to stabilize the LLMs' distribution rather than enhance their core capabilities; 2) \textbf{Incorporating uncertainty representation within the reward model.} Such uncertainty could serve as a crucial link connecting weak reward models with weak-to-strong learning algorithms; 3) \textbf{Efficient construction of high-quality process reward models} that deliver detailed training signals related to the reasoning process \citep{lightman2023let,wang2023math}.

\section{Conclusion, Limitation, and Future Work} We introduce DeepSeekMath, which surpasses all open-source models on the competition-level MATH benchmark and nears the performance of proprietary models. DeepSeekMath~starts from DeepSeek-Coder-v1.5 7B and undergoes continual training over 500B tokens, with a substantial portion of the training data consisting of 120B math tokens obtained from Common Crawl. Our thorough ablation study demonstrates that web pages provide considerable promise for acquiring high-quality mathematical data, whereas arXiv might be less beneficial than initially anticipated. We propose Group Relative Policy Optimization (GRPO), a variant of Proximal Policy Optimization (PPO), which significantly enhances mathematical reasoning abilities while using less memory. Experimental results indicate that GRPO remains effective even when DeepSeekMath-Instruct 7B attains a high benchmark score. Additionally, we offer a unified framework to interpret a range of methods and identify several promising avenues for more efficient reinforcement learning.

While DeepSeekMath~achieves strong results on quantitative reasoning benchmarks, its performance in geometry and theorem-proving tasks is relatively weaker compared to proprietary models. For example, during our preliminary tests, the model struggled with problems involving triangles and ellipses, suggesting potential data selection bias during pre-training and fine-tuning. Moreover, limited by model size, DeepSeekMath~performs worse than GPT-4 in few-shot learning scenarios. GPT-4 can enhance its performance with few-shot examples, whereas DeepSeekMath~shows comparable outcomes in zero-shot and few-shot evaluations. Moving forward, we plan to refine our data selection pipeline to build a higher-quality pre-training corpus. Furthermore, we intend to investigate potential directions (Section \ref{sec:effec_rl}) for more effective reinforcement learning in large language models.

\end{document}